\paragraph{Relation to impurity minimization.}
For a hub state put \(p_i=f_{ih}/m_h\) and append \(p_0=1-\sum_i p_i\) for all unmonitored outcomes. Its contribution in \eqref{eq:allocation} is half its mass times Rao's quadratic diversity \cite{rao}, with all costs involving category zero set to zero. For saturated traffic and unit off-diagonal cost this is the usual weighted Gini impurity up to a factor of two \cite{cicalese}. With slack the enlarged function need not be jointly concave, even for two ports. In coordinates \((p_1,p_2,p_0)\) it is \(p_1p_2\). Its value at the midpoint of \((1/2,1/2,0)\) and \((0,0,1)\) is \(1/16\), below their average \(1/8\). Hence the concave-impurity partition theorem \cite{burshtein} does not apply directly to this augmented simplex. The separate row-concavity argument handles split traffic and allocates the unobserved mass while preserving each measured port total.

\paragraph{A same-observable rank comparison.}
The midpoint contribution of \(B\) to the full matrix is \(G=F\operatorname{diag}(m)^{-1}F^\top\), with entries \(F_{ih}=f_{ih}\). This factorization and its rank bound \(\operatorname{rank}G\le k\) are immediate. For equal traffic \(a\), weighted Cauchy--Schwarz gives \(G\preceq aI\): in each column, \((\sum_i f_{ih}z_i)^2/m_h\le\sum_i f_{ih}z_i^2\). Consequently \(\operatorname{tr}G\le ka\), while \(\sum_{ij}G_{ij}\ge(da)^2/M\). Unit pair weights thus give the elementary rank-based floor
\[
 Y\ge \tfrac12\max\{0,(da)^2/M-ka\}.
\]
At saturation \(M=da\) this recovers \eqref{eq:uniform}. At \(d=4,k=3,M=1,a=0.2\) the rank floor is \(0.02\); dropping the occupancy constraint from \eqref{eq:finite} gives \(0.04\) by \eqref{eq:curv}, and the sharp floor is \(1/15\). An observed score \(0.05\) separates the three-state class only under the last of these tests. This score is itself feasible: with four hub states of mass \(1/4\), interpolate the traffic matrix between \(a I\) and the matrix with all entries \(a/4\), retaining self-loops for unused row mass. The score varies continuously from zero to \(6a^2\). This comparison concerns the specified marginals and scalar. Methods with dwell times, midpoint-labelled triples or a full trajectory have access to other information.

\begin{proposition}[Connected realizations]\label{prop:connected}
Suppose each port has mass strictly greater than \(a_i\), and the unmonitored exterior has positive mass that can be distributed over arbitrarily many states. The minimum in Theorem~\ref{thm:finite} is attained by an irreducible aperiodic kernel. With no unmonitored exterior and \(\sum_i a_i<M\), the same value is the infimum over irreducible aperiodic kernels.
\end{proposition}
\begin{proof}
In the star construction, allocate the exterior mass among \(d\) new states \(e_i\). Join \(e_i\) to port \(i\) and join the exterior states in a path, using small positive symmetric edge masses deducted from endpoint self-loops. Every hub group is connected to this path. Keep a positive self-loop for aperiodicity. Each exterior midpoint still sees at most one port and no port midpoint sees another port. Thus all monitored cross terms and the one-step traffic stay fixed. Without exterior mass, let \(P^*\) be an extremizer on its stars. On the same masses the allocation \(f_{ih}=a_i m_h/M\) connects every hub state to every port with positive self-loops under the strict margins. A convex mixture of its kernel and \(P^*\) preserves all one-step traffic and tends to \(P^*\).
\end{proof}
